\chapter{Pericardiocentesis}

\section*{Overview}
Pericardiocentesis is a procedure that uses a needle and catheter to remove fluid from the pericardial sac around the heart, either for diagnosis or to relieve cardiac tamponade.

\section*{Alternative Names}
Pericardial tap; pericardial fluid drainage

\section*{Principle}
A needle is advanced into the pericardial space, usually with echocardiographic guidance, to aspirate fluid. Drainage reduces intrapericardial pressure and restores cardiac filling in tamponade.

\section*{History}
Pericardiocentesis was described in the 1840s and refined with the introduction of echocardiographic guidance and catheter drainage in the late 20th century.

\section*{Why It Is Done}
Performed urgently in cardiac tamponade, and for diagnostic evaluation of pericardial effusion of uncertain cause, including suspected infection, malignancy, or inflammation.

\section*{Is This a Primary or Secondary Test or Procedure}
Primary emergency procedure for tamponade and a diagnostic tool for pericardial effusion.

\section*{How It Is Performed}
Performed with echocardiographic or fluoroscopic guidance, typically via a subxiphoid or apical approach. The needle enters the pericardial space, fluid is aspirated, and a drain may be left in place.

\section*{Preparation}
Echocardiography to confirm the effusion and guide the approach; emergency equipment is prepared for complications.

\section*{Contraindications and Precautions}
Coagulopathy and a posterior effusion are relative contraindications in the elective setting. Precaution: tamponade is an emergency in which the procedure is performed despite other risks.

\section*{Risks}
Cardiac puncture, pneumothorax, bleeding, arrhythmia, infection, and reaccumulation of fluid.

\section*{Aftercare and Recovery}
The patient is monitored; a drain may remain for days. Fluid is sent for cell count, cytology, culture, and chemistry.

\section*{Outcome and Follow-up}
Findings may indicate infection (purulent or culture-positive), malignancy (positive cytology), uremia, or inflammation. Rapid clinical improvement follows tamponade relief.

\section*{Expected Results}
A small, physiologic amount of clear fluid with no evidence of infection, malignancy, or bleeding.

\section*{Follow-up}
Repeat echocardiography; the underlying cause is treated.

\section*{When to Seek Medical Care}
Follow the procedure team's specific instructions. Seek urgent medical assessment for severe or worsening pain, uncontrolled bleeding, fever or signs of infection, trouble breathing, chest pain, fainting, new weakness or confusion, inability to urinate, or any rapidly worsening symptom. The appropriate warning signs depend on the procedure and the patient's medical condition.

\section*{References}
\begin{enumerate}
  \item MedlinePlus. Pericardiocentesis. U.S. National Library of Medicine; 2025.
  \item Current Medical Diagnosis and Treatment. McGraw-Hill; 2025.
\end{enumerate}

\clearpage
